\documentclass[preprint,12pt]{elsarticle}

\usepackage[utf8]{inputenc}
\usepackage[T1]{fontenc}
\usepackage{lmodern}
\usepackage{amsmath}
\usepackage{booktabs}
\usepackage{graphicx}
\usepackage{array}
\usepackage{url}
\usepackage{xurl}
\usepackage{microtype}

\journal{Environmental Modelling \& Software}

\renewcommand{\thesection}{S\arabic{section}}
\renewcommand{\thetable}{S\arabic{table}}
\setcounter{secnumdepth}{2}

\begin{document}
\setlength{\emergencystretch}{3em}

\begin{frontmatter}
\title{Supporting Information for: Discard or reconstruct? A leakage-safe benchmark of gap handling for annual heat extremes and their trends in river-temperature archives}
\end{frontmatter}

\noindent This Supporting Information contains eleven tables that document the station roster and trend signal, the archive-level direction counts, the scope of the uncertainty summaries, the H7 refit robustness check, the open-source and flow diagnostics, the definitions, the software release, and the data-access terms. All values are taken from saved result files; no new model fit or data download was performed for these tables.

\section{Station, group, and reference-year summary}
Table~\ref{tab:S1} summarizes the 41 forward stations by region and connected group. The cohort contains 24 Bavarian stations in eight groups, 10 Swiss stations in two groups, and seven USGS stations in seven groups. The median reference-year set contains 14 complete years (minimum 7, maximum 16), and 24 stations have at least one internal gap year inside their reference span.

\begin{table}[t]
\centering
\small
\caption{Station, connected-group, and reference-year summary by region. Complete-record years are the years in each station's saved reference-year set.}
\label{tab:S1}
\resizebox{\linewidth}{!}{%
\begin{tabular}{lrrrrrr}
\toprule
Region & Stations & Groups & \shortstack{Median complete-\\record years (range)} & \shortstack{Median first--\\last year} & \shortstack{Median\\span (years)} & \shortstack{Stations with\\internal gaps} \\
\midrule
FOEN (Switzerland) & 10 & 2 & 16 (15--16) & 2010--2025 & 16 & 1 \\
GKD Bayern (Bavaria) & 24 & 8 & 13 (7--16) & 2010--2025 & 15 & 18 \\
USGS HUC03 & 1 & 1 & 14 (14--14) & 2011--2025 & 15 & 1 \\
USGS HUC04 & 1 & 1 & 11 (11--11) & 2010--2024 & 15 & 1 \\
USGS HUC10 & 2 & 2 & 14 (12--16) & 2010--2025 & 16 & 1 \\
USGS HUC17 & 3 & 3 & 12 (10--14) & 2010--2024 & 14 & 2 \\
\midrule
All forward stations & 41 & 17 & 14 (7--16) & 2010--2025 & 15 & 24 \\
\bottomrule
\end{tabular}}
\par\smallskip\footnotesize Notes: A station-year is one station and one calendar year. A connected group links stations by river network or record family and is the cluster unit. The reference-year set contains the complete years available for a station under the design and need not be contiguous. Span is last year minus first year plus one. Internal gaps count stations with at least one missing year inside the reference span. Source: \path{results/wrr_final_revision_20261031/reference_year_sets.csv} and \path{results/wrr_final_revision_20261031/retention_choice_counts.csv}.
\end{table}


\section{Full station-level trend metrics}
Table~\ref{tab:S2} reports the complete-record H30 trend metrics for all 41 forward stations. The OLS slope and its classical standard error are computed on each station's saved reference-year set. The Sen slope is the median of pairwise slopes. The Mann--Kendall $p$-value includes a tie correction in the variance of $S$ and a continuity correction, but no autocorrelation correction. The H30 interannual SD describes the year-to-year variability of the complete-record annual maximum.

\begin{table}[p]
\centering
\footnotesize
\caption{Full station-level H30 trend metrics for the 41 forward stations. OLS slope, OLS SE, and Sen slope are in $^{\circ}$C/10a; H30 interannual SD is in $^{\circ}$C. OLS SE is classical and excludes autocorrelation. Mann--Kendall $p$ includes a tie correction in $\operatorname{Var}(S)$ but no autocorrelation correction.}
\label{tab:S2}
\renewcommand{\arraystretch}{0.96}
\resizebox{\linewidth}{!}{%
\begin{tabular}{lllrrrrrrr}
\toprule
Station & Region & Group & $n$ & OLS & OLS & OLS & Sen & MK & H30 \\
 & & & & slope & SE & $p$ & slope & $p$ & SD \\
\midrule
\texttt{foen|2016} & FOEN & g0029 & 16 & 0.850 & 0.644 & 0.208 & 0.863 & 0.065 & 1.217 \\
\texttt{foen|2019} & FOEN & g0029 & 16 & 0.496 & 0.180 & 0.015 & 0.557 & 0.027 & 0.398 \\
\texttt{foen|2029} & FOEN & g0029 & 16 & 0.757 & 0.645 & 0.260 & 0.816 & 0.079 & 1.204 \\
\texttt{foen|2030} & FOEN & g0029 & 16 & 0.960 & 0.543 & 0.099 & 0.633 & 0.043 & 1.070 \\
\texttt{foen|2085} & FOEN & g0029 & 16 & 0.349 & 0.619 & 0.582 & 0.331 & 0.499 & 1.115 \\
\texttt{foen|2113} & FOEN & g0029 & 15 & 0.780 & 0.790 & 0.342 & 0.728 & 0.138 & 1.385 \\
\texttt{foen|2126} & FOEN & g0030 & 16 & 0.423 & 0.541 & 0.447 & 0.261 & 0.893 & 0.984 \\
\texttt{foen|2135} & FOEN & g0029 & 16 & 0.841 & 0.549 & 0.148 & 0.723 & 0.115 & 1.057 \\
\texttt{foen|2205} & FOEN & g0029 & 16 & 0.906 & 0.648 & 0.184 & 0.972 & 0.053 & 1.232 \\
\texttt{foen|2457} & FOEN & g0029 & 16 & 1.248 & 0.279 & $<$0.001 & 1.356 & $<$0.001 & 0.774 \\
\texttt{gkd\_bayern|10026293} & GKD & g0032 & 13 & 0.823 & 0.712 & 0.273 & 0.554 & 0.246 & 1.148 \\
\texttt{gkd\_bayern|10039802} & GKD & g0032 & 12 & 0.712 & 0.808 & 0.399 & 0.650 & 0.304 & 1.199 \\
\texttt{gkd\_bayern|10046105} & GKD & g0032 & 15 & 1.235 & 0.696 & 0.099 & 1.261 & 0.113 & 1.251 \\
\texttt{gkd\_bayern|10056302} & GKD & g0032 & 16 & 0.643 & 0.661 & 0.348 & 0.513 & 0.344 & 1.217 \\
\texttt{gkd\_bayern|10074009} & GKD & g0032 & 13 & 0.465 & 0.908 & 0.619 & 1.097 & 0.360 & 1.326 \\
\texttt{gkd\_bayern|10081004} & GKD & g0032 & 14 & 0.745 & 0.940 & 0.443 & 0.585 & 0.443 & 1.397 \\
\texttt{gkd\_bayern|10089006} & GKD & g0032 & 15 & 0.568 & 0.853 & 0.517 & 0.575 & 0.621 & 1.398 \\
\texttt{gkd\_bayern|10091008} & GKD & g0032 & 14 & 1.081 & 1.037 & 0.318 & 1.003 & 0.274 & 1.710 \\
\texttt{gkd\_bayern|10098007} & GKD & g0032 & 14 & $-0.411$ & 1.066 & 0.706 & $-0.767$ & 0.511 & 1.555 \\
\texttt{gkd\_bayern|11401009} & GKD & g0034 & 12 & 0.539 & 0.405 & 0.212 & 0.718 & 0.304 & 0.656 \\
\texttt{gkd\_bayern|11402001} & GKD & g0034 & 14 & 1.204 & 0.688 & 0.106 & 1.642 & 0.101 & 1.239 \\
\texttt{gkd\_bayern|11405000} & GKD & g0034 & 7 & 0.544 & 0.508 & 0.333 & 0.233 & 0.368 & 0.689 \\
\texttt{gkd\_bayern|16003003} & GKD & g0036 & 13 & 0.867 & 0.453 & 0.082 & 0.813 & 0.127 & 0.818 \\
\texttt{gkd\_bayern|16005701} & GKD & g0036 & 9 & 1.542 & 0.953 & 0.150 & 1.581 & 0.251 & 1.433 \\
\texttt{gkd\_bayern|16007004} & GKD & g0036 & 16 & 0.819 & 0.450 & 0.090 & 0.860 & 0.137 & 0.891 \\
\texttt{gkd\_bayern|17207508} & GKD & g0039 & 9 & 0.930 & 0.529 & 0.122 & 0.549 & 0.348 & 0.706 \\
\texttt{gkd\_bayern|18001508} & GKD & g0035 & 8 & 0.568 & 0.673 & 0.431 & 0.669 & 0.266 & 0.509 \\
\texttt{gkd\_bayern|18008008} & GKD & g0035 & 10 & 1.037 & 0.775 & 0.218 & 1.004 & 0.107 & 1.091 \\
\texttt{gkd\_bayern|18403002} & GKD & g0031 & 14 & 1.078 & 0.584 & 0.090 & 1.167 & 0.101 & 1.011 \\
\texttt{gkd\_bayern|24010004} & GKD & g0037 & 12 & 0.757 & 0.590 & 0.229 & 0.494 & 0.304 & 1.032 \\
\texttt{gkd\_bayern|24022003} & GKD & g0037 & 10 & 0.564 & 0.603 & 0.377 & 0.407 & 0.721 & 0.690 \\
\texttt{gkd\_bayern|24042000} & GKD & g0037 & 12 & $-0.026$ & 0.584 & 0.965 & $-0.061$ & 1.000 & 0.791 \\
\texttt{gkd\_bayern|24064003} & GKD & g0037 & 14 & 0.754 & 0.507 & 0.162 & 0.743 & 0.228 & 0.946 \\
\texttt{gkd\_bayern|24405002} & GKD & g0033 & 9 & 0.795 & 0.653 & 0.263 & 0.883 & 0.348 & 0.952 \\
\texttt{usgs|02172020} & HUC03 & g0181 & 14 & 0.398 & 0.270 & 0.166 & 0.467 & 0.125 & 0.463 \\
\texttt{usgs|04137030} & HUC04 & g0166 & 11 & $-0.609$ & 0.317 & 0.086 & $-0.750$ & 0.087 & 0.529 \\
\texttt{usgs|06036905} & HUC10 & g0053 & 16 & 0.701 & 0.415 & 0.113 & 0.438 & 0.224 & 0.811 \\
\texttt{usgs|06054500} & HUC10 & g0143 & 12 & 0.300 & 0.531 & 0.584 & $-0.243$ & 0.945 & 0.791 \\
\texttt{usgs|13342500} & HUC17 & g0054 & 12 & 0.273 & 0.439 & 0.548 & 0.000 & 1.000 & 0.691 \\
\texttt{usgs|14166000} & HUC17 & g0057 & 14 & 0.426 & 0.487 & 0.399 & 0.467 & 0.381 & 0.727 \\
\texttt{usgs|14211500} & HUC17 & g0059 & 10 & 1.776 & 0.468 & 0.005 & 1.638 & 0.004 & 0.874 \\
\bottomrule
\end{tabular}}
\par\smallskip\footnotesize Notes: $n$ is the number of reference years. $p$-values below 0.001 are printed as $<$0.001. Values are rounded for display; the saved file retains full precision. Source: \path{results/ems_revision_20261102/trend_signal.csv}. Region labels group the USGS stations into HUC03, HUC04, HUC10, and HUC17. Station keys retain the provider prefix.
\end{table}


\section{Archive-level direction counts}
Table~\ref{tab:S3} gives the within-archive direction of the all-year comparison for both gap families and for the same-count CDF and $S_{xx}$ calendar controls. The all-year rows are the primary comparison. Numerical equality uses an arithmetic $10^{-12}\,^{\circ}$C per decade tolerance and is not practical equivalence. Archives are station/mask combinations, not independent rivers. The uniform-exact calendar control is not tabulated because its archive-level estimate errors are not saved; its cohort RMSE is 0.449 (controlled) and 0.135 (calendar) $^{\circ}$C/10a, with 55 and 8 sign flips.

\begin{table}[p]
\centering
\footnotesize
\caption{Archive-level direction of the all-year slope-error comparison. Counts give archives in which each scheme's absolute finite-record slope error is smaller than, larger than, or numerically equal to warm-season screening; the all-year rows are the primary comparison and the CDF and $S_{xx}$ rows use the same-count controls. Median change is the within-archive absolute-error change in $^{\circ}$C per decade. Numerical equality uses an arithmetic $10^{-12}\,^{\circ}$C per decade tolerance, which is not practical equivalence. The controlled 30--90-day gap and empirical-calendar families are separate cohorts, and archives are station/mask combinations rather than independent rivers; a changed year set means that the selected years differ from those retained under warm-season screening.}
\label{tab:S3}
\resizebox{\linewidth}{!}{%
\begin{tabular}{llrrrrrr}
\toprule
Family & Scheme & Subset & $n$ & Smaller abs. error & Larger abs. error & Numerically equal & Median change \\
\midrule
Controlled 30--90 day & All years & All archives & 805 & 630 & 169 & 6 & $-0.1395$ \\
Controlled 30--90 day & All years & Changed year set & 800 & 630 & 169 & 1 & $-0.1400$ \\
Controlled 30--90 day & CDF & All archives & 805 & 197 & 170 & 438 & 0.0000 \\
Controlled 30--90 day & CDF & Changed year set & 367 & 197 & 170 & 0 & $-0.0196$ \\
Controlled 30--90 day & $S_{xx}$ & All archives & 805 & 218 & 158 & 429 & 0.0000 \\
Controlled 30--90 day & $S_{xx}$ & Changed year set & 376 & 218 & 158 & 0 & $-0.0421$ \\
\addlinespace
Empirical calendar & All years & All archives & 818 & 231 & 57 & 530 & 0.0000 \\
Empirical calendar & All years & Changed year set & 288 & 231 & 57 & 0 & $-0.0578$ \\
Empirical calendar & CDF & All archives & 818 & 144 & 91 & 583 & 0.0000 \\
Empirical calendar & CDF & Changed year set & 235 & 144 & 91 & 0 & $-0.0256$ \\
Empirical calendar & $S_{xx}$ & All archives & 818 & 151 & 83 & 584 & 0.0000 \\
Empirical calendar & $S_{xx}$ & Changed year set & 234 & 151 & 83 & 0 & $-0.0490$ \\
\bottomrule
\end{tabular}}
\par\smallskip\footnotesize Notes: ``Smaller abs. error'' and ``Larger abs. error'' count archives relative to warm-season screening at the same station and mask. The all-year comparison is the primary row; the CDF and $S_{xx}$ rows are same-count calendar controls. Median change is the median within-archive change in absolute slope error, not a cohort RMSE. The \texttt{uniform\_exact} scheme is not shown as a row because its archive-level estimate errors are not saved for this cohort; its saved cohort RMSE is 0.449 (controlled) and 0.135 (calendar) $^{\circ}$C/10a, with 55 and 8 sign flips. Source: \path{results/wrr_comprehensive_review_20261102/archive_level_comparisons.csv}; the uniform-exact cohort values are from \path{results/ems_revision_20261102/trend_signal.csv}.
\end{table}


\section{Scope of the uncertainty summaries}
Table~\ref{tab:S4} states what each of the four uncertainty summaries describes and what it does not establish. The four summaries are annual H30 90\% intervals, finite-record slope 95\% intervals, paired RMSE-difference 95\% group intervals, and three-source-seed or random-path ranges. None is a general coverage guarantee or a climate-process confidence interval.

\begin{table}[p]
\centering
\footnotesize
\caption{Scope of four uncertainty summaries. Annual H30 90\% intervals describe the annual predictive distribution under fixed methods and information. Finite-record slope 95\% intervals propagate that annual distribution and the specified joint construction to the slope. Paired RMSE-difference 95\% group intervals describe comparison variability conditional on the saved annual predictions, the predefined connected groups and the rules compared. Three-source-seed or random-path ranges describe numerical simulation variation at a fixed budget. None establishes a general coverage guarantee, a climate-process confidence interval, or uncertainty from retraining source models and selecting the rule. At 4,096 samples, the nominal 95\% trend intervals remain 87.4--89.9\%, below nominal 95\%, as a saved diagnostic rather than a new experiment.}
\label{tab:S4}
\setlength{\tabcolsep}{4pt}
\begin{tabular}{>{\raggedright\arraybackslash}p{0.16\textwidth}>{\raggedright\arraybackslash}p{0.35\textwidth}>{\raggedright\arraybackslash}p{0.39\textwidth}}
\toprule
Summary & What it describes & What it does not establish \\
\midrule
Annual H30 90\% intervals & Annual predictive distribution under fixed methods and information; empirical coverage of the annual target & Not evidence that every region or reconstructor attains 90\%; endpoint comparisons use an arithmetic tolerance. \\
\addlinespace
Finite-record slope 95\% intervals & Propagates the annual distribution and the specified joint construction to the slope on the reference-year set & Not a 95\% confidence guarantee for a climate-process true trend, another reference-year set or a complete calendar. \\
\addlinespace
Paired RMSE-difference 95\% group intervals & Comparison variability conditional on saved annual predictions, predefined connected groups and the rules compared & Not retraining or rule-selection uncertainty, multiplicity across sensitivities, or new climate samples. \\
\addlinespace
Three-source-seed / random-path ranges & Numerical or simulation variation at a stated budget with fixed data and rules & Not a confidence interval across independent hydrological samples; more paths do not establish convergence. \\
\bottomrule
\end{tabular}
\par\smallskip\footnotesize Notes: The 87.4--89.9\% range gives the minimum lower and maximum upper bounds of the 4,096-draw coverage intervals across the saved methods and constructions, not point estimates; the corresponding point estimates span 87.5--89.7\%. It is a saved diagnostic for the nominal 95\% trend intervals, not a new experiment or a general coverage guarantee. Source: \path{results/wrr_adversarial_review_20261027/rank_convergence/rank_convergence_summary.csv}.
\end{table}


\section{H7 refit robustness}
Table~\ref{tab:S5} reports the H7 refit. The H7 target is the annual maximum of seven consecutive daily means, not the H30 target used in the main analysis. The refit used 20 seeds and 540 outer-region fits. Panel A gives the danger-scope MAE, CRPS, tail CRPS, and 90\% coverage for the four estimators. Panel B gives the direction-stability comparison of D2 relative to the stochastic comparator (calibrated stochastic regression): the median paired MAE reduction is 2.94\% in the crossed family and 4.85\% in the calendar family, and CRPS and tail CRPS improve in 20 of 20 seeds in both families. Panel C gives trend-target missingness-component coverage: D2 multiple imputation has median coverage 81.2\% under MCAR and 80.7\% under the strong-heat scenario, but 4.5\% under the non-air-trend scenario. The near-81\% value therefore describes the MCAR and strong-heat trend scenarios, not the non-air-trend scenario.

\begin{table}[p]
\centering
\footnotesize
\caption{H7 refit robustness. Panel A reports native H7 source refits over 20 seeds and 540 outer-region fits; seven-day peak supports, masks, history summaries, conditional functionals, MOVE augmentation, annual models, and disjoint calibration were recomputed. Panel B reports the direction-stability check for D2 (distribution regression) relative to the stochastic comparator (calibrated stochastic regression). Panel C reports trend-target coverage of the missingness component for the multiple-imputation pipelines. These are H7 results, not the historical H30 endpoint.}
\label{tab:S5}

\textit{Panel A: H7 danger-scope refit scores (20 seeds).}\par\smallskip
\resizebox{\linewidth}{!}{%
\begin{tabular}{llrrrrr}
\toprule
Family & Estimator & Seeds & MAE ($^{\circ}$C) & CRPS ($^{\circ}$C) & Tail CRPS ($^{\circ}$C) & Coverage 90\% \\
\midrule
Calendar & Distribution regression (D2) & 20 & 0.494 & 0.365 & 0.065 & 0.874 \\
Calendar & Functional debiasing & 20 & 0.522 & -- & -- & -- \\
Calendar & Daily plug-in & 20 & 0.638 & -- & -- & -- \\
Calendar & Calibrated stochastic & 20 & 0.516 & 0.380 & 0.068 & 0.858 \\
Crossed & Distribution regression (D2) & 20 & 0.348 & 0.257 & 0.046 & 0.892 \\
Crossed & Functional debiasing & 20 & 0.375 & -- & -- & -- \\
Crossed & Daily plug-in & 20 & 0.423 & -- & -- & -- \\
Crossed & Calibrated stochastic & 20 & 0.360 & 0.265 & 0.048 & 0.893 \\
\bottomrule
\end{tabular}}

\par\medskip
\textit{Panel B: D2 relative to the stochastic comparator (calibrated stochastic regression).}\par\smallskip
\resizebox{\linewidth}{!}{%
\begin{tabular}{llrrr}
\toprule
Family & Comparison & \shortstack{Median paired\\MAE reduction (\%)} & \shortstack{CRPS gain\\positive seeds} & \shortstack{Tail CRPS gain\\positive seeds} \\
\midrule
Crossed & D2 vs calibrated stochastic & 2.94 & 20/20 & 20/20 \\
Calendar & D2 vs calibrated stochastic & 4.85 & 20/20 & 20/20 \\
\bottomrule
\end{tabular}}

\par\medskip
\textit{Panel C: Trend-target missingness-component coverage (median, seed range).}\par\smallskip
\resizebox{\linewidth}{!}{%
\begin{tabular}{lrrr}
\toprule
Scenario & D2 multiple imputation & Calibrated stochastic MI & Stochastic regression MI \\
\midrule
MCAR & 81.2\% (79.1--86.2) & 78.4\% (75.7--79.7) & 75.6\% (73.9--76.7) \\
Strong heat & 80.7\% (78.9--87.4) & 78.5\% (75.5--81.3) & 76.0\% (73.8--78.4) \\
Non-air trend & 4.5\% (3.8--5.6) & 3.9\% (2.9--4.5) & 3.7\% (2.7--4.2) \\
\bottomrule
\end{tabular}}
\par\smallskip\footnotesize Notes: Panel A values are medians over 20 seeds in the dangerous decision-time stratum. Coverage is the empirical coverage of nominal 90\% annual H7 intervals. Panel B ``median paired MAE reduction'' is the median over saved station-level paired relative reductions in MAE for D2 relative to the calibrated stochastic comparator across the 20 seeds; positive CRPS and tail-CRPS gains mean lower D2 scores. Panel C covers the trend target, not the annual H7 target; values are medians with minimum--maximum seed ranges. H7 is a different target and is not used as confirmation of the H30 results. Source: \path{results/wrr_review_completion_20261024/evidence/H7_refit_scorecard_summary.csv}, \path{results/wrr_review_completion_20261024/evidence/H7_refit_direction_stability.csv}, \path{results/wrr_review_completion_20261024/target_windows/h7/full/trends_per_seed.csv}, and \path{results/wrr_review_completion_20261024/target_windows/h7/full/completion.json}.
\end{table}


\section{Open-source candidate-provider sensitivity}
Table~\ref{tab:S6} reports the candidate-provider sensitivity. Replacing the restricted source providers with seven candidate providers changes the crossed/forward D2 trend RMSE from 0.164 to 0.158\,$^{\circ}$C/10a and the classical trend RMSE from 0.147 to 0.156\,$^{\circ}$C/10a over 738 four-way common pairs in 17 groups. D2 retention falls from 93.9\% to 93.7\% on 807 factor cases, and nominal 95\% D2 interval coverage falls from 88.0\% to 87.6\% on 820 coverage cases. The ranking of the two approaches is unchanged. This is a candidate open reproducibility route, not a legal clearance.

\begin{table}[t]
\centering
\small
\caption{Open-source candidate-provider sensitivity on the crossed/forward cohort. The candidate route refits the source models on seven candidate providers with recorded reuse paths while retaining the target masks, daily features, and plug-in baselines. Retention and coverage are D2 quantities.}
\label{tab:S6}
\resizebox{\linewidth}{!}{%
\begin{tabular}{lrr}
\toprule
Metric & Original source & Candidate source \\
\midrule
D2 trend RMSE ($^{\circ}$C/10a) & 0.164 & 0.158 \\
Classical trend RMSE ($^{\circ}$C/10a) & 0.147 & 0.156 \\
D2 retention in the symmetric MSE decomposition (\%) & 93.9 & 93.7 \\
Nominal 95\% D2 interval coverage (\%) & 88.0 & 87.6 \\
Common four-way pairs & 738 & 738 \\
Connected groups & 17 & 17 \\
Factor cases for retention & 807 & 807 \\
Coverage cases & 820 & 820 \\
\bottomrule
\end{tabular}}
\par\smallskip\footnotesize Notes: The candidate route leaves the ranking of D2 and the classical comparator unchanged. It is a candidate open reproducibility route, not a legal clearance: LUBW rights and the provenance of neighbor records still require confirmation. Candidate-source tables are descriptive, repetitions are not independent climates, and the filtered inputs remain processed study data rather than provider-native records. Source: \path{results/ems_revision_20261102/flow_diagnostic.json}, extracted from \path{paper/wrr_tr/open_source_values.tex}.
\end{table}


\section{Zero-flow flow-proxy logistic results}
Table~\ref{tab:S7} reports the station-fixed conditional-logit associations between a warm-season-incomplete station-year and within-station standardized summer air temperature, within-station standardized summer log flow, and a linear trend. The flow-adjusted odds ratios are 0.594 [0.409, 0.813] for the recorded-zero cohort and 0.474 [0.265, 0.676] for the at-least-30-zero-day cohort. The air-temperature odds ratios are 0.993 [0.659, 1.305] and 0.855 [0.580, 1.275]. These are associations in a historical proxy cohort, not causal effects or a missing-at-random test.

\begin{table}[t]
\centering
\small
\caption{Zero-flow flow-proxy conditional-logit results. The response is a warm-season-incomplete station-year (the hidden-peak condition). Air and flow covariates are within-station standardized summer means; the model also includes a linear trend per decade. Odds ratios are shown with 95\% confidence intervals.}
\label{tab:S7}
\resizebox{\linewidth}{!}{%
\begin{tabular}{lrrrrrr}
\toprule
Cohort & $n$ & Events & Stations & Groups & $z_{\mathrm{air}}$ OR (95\% CI) & $z_{\mathrm{flow}}$ OR (95\% CI) \\
\midrule
Recorded zero flow & 282 & 129 & 21 & 17 & 0.993 (0.659--1.305) & 0.594 (0.409--0.813) \\
At least 30 zero days & 215 & 104 & 16 & 14 & 0.855 (0.580--1.275) & 0.474 (0.265--0.676) \\
Fewer than 30 zero days & 4,417 & 1,814 & 271 & 103 & 1.024 (0.944--1.122) & 0.969 (0.897--1.046) \\
No zero in available flow & 4,350 & 1,789 & 266 & 103 & 1.020 (0.937--1.125) & 0.968 (0.895--1.050) \\
\bottomrule
\end{tabular}}
\par\smallskip\footnotesize Notes: The underlying USGS missingness sample contains 6,479 station-years at 336 physical stations; 21 stations in 17 groups recorded at least one zero-flow day and 16 stations in 14 groups recorded at least 30. Confidence intervals use 400 bootstrap draws. Lower hidden-peak odds is an association in this proxy cohort, not evidence that low flow protects against underestimation. Zero flow is a historical proxy; no recorded zero does not prove perennial flow, and a zero does not identify an intermittent river, sensor exposure, or a dry bed. Source: \path{results/wrr_review_completion_20261024/zero_flow_missingness/conditional_logit_summary.csv} and \path{results/wrr_review_completion_20261024/zero_flow_missingness/cohort_counts.csv}.
\end{table}


\section{Flow-comparator summary}
Table~\ref{tab:S8} reports availability and mean absolute error for the opened flow-comparator methods. Historical D2 has full availability and the lowest MAE in both scopes. Donor-dependent methods are available on 81.7\% of crossed/all eligible cases and do not lower paired MAE. The paired comparisons use 66,099 common cases in the all scope and 33,609 in the danger scope, covering 210 sites. The full cohort contains 255 sites, 8,999 complete climatic years, 134,925 cases, and 17 HUC2 region groups.

\begin{table}[t]
\centering
\small
\caption{Opened flow-comparator scorecard for the crossed family. Availability is the fraction of eligible cases with a valid estimate; MAE is mean absolute error in log-flow units computed on each estimator's valid cases. The ``all'' scope contains 80,931 eligible cases and the ``danger'' scope 41,132; paired log-MAE comparisons use 66,099 and 33,609 common cases, respectively.}
\label{tab:S8}
\resizebox{\linewidth}{!}{%
\begin{tabular}{lrrrr}
\toprule
Estimator & \shortstack{Availability,\\all} & \shortstack{MAE,\\all} & \shortstack{Availability,\\danger} & \shortstack{MAE,\\danger} \\
\midrule
Historical D2 & 1.000 & 0.0500 & 1.000 & 0.0743 \\
Historical multidonor Ridge & 1.000 & 0.0564 & 1.000 & 0.0829 \\
Single-donor Ridge & 0.817 & 0.0602 & 0.817 & 0.0896 \\
Observed days & 1.000 & 0.0646 & 1.000 & 0.1004 \\
Stochastic AR(1) median & 0.817 & 0.0716 & 0.817 & 0.1038 \\
MOVE.1 & 0.817 & 0.0734 & 0.817 & 0.1062 \\
\bottomrule
\end{tabular}}
\par\smallskip\footnotesize Notes: The tabulated MAE is log-MAE (mean absolute error in log-flow units) on each estimator's valid cases: all eligible cases for historical D2, historical multidonor Ridge, and observed days; the donor-available subset for single-donor Ridge, stochastic AR(1), and MOVE.1 (availability 0.817, i.e., 66,099 of 80,931 eligible cases in the all scope and 33,609 of 41,132 in the danger scope). Paired log-MAE comparisons are restricted to these common cases and cover 210 sites; the main-text paired basis yields D2 MAE 0.0495 (all) and 0.0745 (danger), slightly different from the per-estimator values tabulated here. The full cohort contains 255 sites, 8,999 complete climatic years, 134,925 cases, and 17 HUC2 region groups. The 1.000 availability rows are historical D2, historical multidonor Ridge, and observed days; the donor-dependent rows do not lower paired MAE. The experiment is an opened positive-flow perennial subset, not a new blind transfer or a 20-seed stability claim. Paired differences are stored as left and right MAE on common cases; no derived subtraction is saved. Source: \path{results/wrr_review_completion_20261024/flow_comparators/scorecard.csv} and \path{results/wrr_review_completion_20261024/flow_comparators/paired_effects.csv}; cohort counts are from \path{results/wrr_review_completion_20261024/flow_comparators/status.json} and \path{results/wrr_review_completion_20261024/flow_comparators/design.json}.
\end{table}


\section{Definitions and symbols}
Table~\ref{tab:S9} collects the terms and symbols used in the main text and in this Supporting Information. The decomposition $E=S+R$ is an arithmetic decomposition of the finite-record slope error, not a causal variance decomposition. The observed lower bound $L_y$ is a property of the reported series, not an annual estimate.

\begin{table}[p]
\centering
\footnotesize
\renewcommand{\arraystretch}{0.95}
\caption{Definitions and symbols used in the main text and this Supporting Information.}
\label{tab:S9}
\begin{tabular}{>{\raggedright\arraybackslash}p{0.20\textwidth}>{\raggedright\arraybackslash}p{0.73\textwidth}}
\toprule
Term or symbol & Definition \\
\midrule
Station-year & One station and one calendar year. \\
Archive & One station $\times$ one source repetition $\times$ one gap setting or empirical calendar. The controlled-gap cohort contains 805 analyzed archives; the empirical-calendar cohort contains 818. \\
Connected group & A set of stations linked by the same river network or record family. It is the cluster unit for group-disjoint calibration and resampling. \\
Reference-year set $Y_i$ & The complete years available for station $i$ under the evaluation design. It need not be contiguous, and years outside it are not imputed as part of the target. The target is the finite-record slope on $Y_i$. \\
H30 & Annual maximum of the 30-day running mean water temperature in year $y$. It is a sustained-heat summary, not an ecological standard. \\
H7 & Annual maximum 7-day mean water temperature. It is a different target, used only as a robustness check and not as confirmation of the H30 results. \\
$L_y$ & Largest fully observed 30-day mean in an incomplete year $y$. It satisfies $\mathrm{H30}_y\ge L_y$ for the reported series and is an observed lower bound, not an annual estimate. \\
SE & Classical OLS standard error of the finite-record slope on the reference years. It excludes autocorrelation and is a lower bound on the trend's sampling uncertainty. \\
$S$, $R$, $E$ & With $\beta_F$ the complete-record slope on $Y_i$, $\beta_S$ the complete-record slope on the selected year set, and $\beta_{SE}$ the estimated slope on that set: $S=\beta_S-\beta_F$, $R=\beta_{SE}-\beta_S$, and $E=\beta_{SE}-\beta_F=S+R$. The squared-error identity includes the cross term $2\,\operatorname{mean}(SR)$; the terms are an arithmetic decomposition, not causal variance shares. \\
Controlled gap & A contiguous 30-, 60-, or 90-day mask applied jointly to every year of an archive, with at least half overlap of the training-period climatological peak. This is a severe stress, not an estimate of outage frequency. \\
Empirical calendar & A year-specific gap calendar copied from the observed archive. It preserves each year's gap length and seasonal position but not the exact missing days, so it is closer to observed annual longest-gap conditions than the controlled stress. \\
Warm-season rule & A station-specific 90\% coverage rule requiring coverage in the five consecutive warm months. The 95\% annual-coverage rule is a secondary comparator. \\
\bottomrule
\end{tabular}
\par\smallskip\footnotesize Notes: Definitions follow the main text and this Supporting Information; the 805 controlled-gap and 818 empirical-calendar archive counts are from \path{results/ems_revision_20261102/trend_signal.csv}.
\end{table}


\section{Software and reproducibility}
Table~\ref{tab:S10} summarizes the \texttt{gapbench} 1.2.0 release. The package exposes six public functions, passes 13 tests, includes a synthetic quickstart, and provides a Docker smoke environment that imports LightGBM and runs the quickstart without restricted observations. The version is consistent across the package metadata and citation file. An archival DOI and a v1.2.0 release tag are pending. The D2 model and the conditional AR(1) bridge remain in the research compendium.

\begin{table}[p]
\centering
\footnotesize
\renewcommand{\arraystretch}{0.95}
\caption{Software and reproducibility summary for \texttt{gapbench} 1.2.0. The released package contains the reusable primitive layer; fitted D2 and AR(1) bridge models remain in the research compendium.}
\label{tab:S10}
\begin{tabular}{>{\raggedright\arraybackslash}p{0.22\textwidth}>{\raggedright\arraybackslash}p{0.71\textwidth}}
\toprule
Item & Status \\
\midrule
Package & \texttt{gapbench} 1.2.0; Python $\ge$ 3.10; NumPy $\ge$ 1.24; pandas $\ge$ 2.0; MIT license for code. \\
Public API & Six functions exported from the package root: \texttt{observed\_window\_bound}, \texttt{finite\_record\_slope}, \texttt{select\_years}, \texttt{decompose\_slope\_error}, \texttt{group\_splits}, and \texttt{mask\_archive}. \\
Tests & 13 tests. Clean test command: \texttt{pip install -e "gapbench[test]"} followed by \texttt{python -m pytest -q gapbench/tests}; all 13 pass. \\
Quickstart & Synthetic 20-year daily series with four gaps. The example measures the best fully observed 30-day window, selects 18 years under both calendar rules while retaining all 16 complete years, and decomposes the slope error. It runs in under one second. \\
Docker & Python 3.11 slim image; installs the review and development extras, imports LightGBM, and runs the quickstart. The container contains no restricted observations. \\
Version consistency & \texttt{pyproject.toml}, \texttt{CITATION.cff}, and the preferred citation report version 1.2.0; the consistency check reports PASS with a warning that the historical v1.0.0 git tag does not describe version 1.2.0. \\
Archival DOI & No Zenodo deposit yet. GitHub release v1.2.0 marks this version; the archival DOI remains pending. \\
Reviewer access & The public repository requires no credentials. Restricted LUBW observations are not redistributed; their access route remains pending provider agreement and editorial/legal confirmation. \\
Program size and replay memory & 24 files, 1{,}012 lines of Python, 264\,KiB uncompressed (23\,KiB zipped). Measured on one CPU core: quickstart 0.13\,s wall clock and 35\,MiB peak resident set; full packaged replay of all 1{,}623 archive-level masks (8{,}115 strategy rows) 9.7\,s and 289\,MiB peak resident set. \\
Research compendium & The D2 quantile model and the conditional AR(1) bridge are not part of \texttt{gapbench}. This separation keeps the published primitives stable while the fitted models depend on restricted data. \\
\bottomrule
\end{tabular}
\par\smallskip\footnotesize Notes: The MIT license applies only to original software and project-authored documentation. It does not relicense observational data or upstream provider products. Source: \path{gapbench/pyproject.toml}, \path{CITATION.cff}, \path{Dockerfile}, and the test suite under \path{gapbench/tests}.
\end{table}


\section{Data and access summary}
Table~\ref{tab:S11} summarizes the data providers and access terms. FOEN, GKD Bayern, USGS NWIS, and NASA POWER are public sources with provider-specific attribution requirements. LUBW observations are restricted; reviewer access remains pending provider agreement and editorial/legal confirmation. The MIT license applies to the \texttt{gapbench} code and does not relicense observational data. Processed annual tables are distributed in the public repository; restricted daily observations are not. Archival DOIs for the code release and processed tables are pending.

\begin{table}[p]
\centering
\footnotesize
\caption{Data and access summary. Provider terms govern access, use, and redistribution. The MIT license applies to code, not to observational data.}
\label{tab:S11}
\setlength{\tabcolsep}{4pt}
\begin{tabular}{>{\raggedright\arraybackslash}p{0.15\textwidth}>{\raggedright\arraybackslash}p{0.18\textwidth}>{\raggedright\arraybackslash}p{0.35\textwidth}>{\raggedright\arraybackslash}p{0.23\textwidth}}
\toprule
Provider & Data & Access and terms & Distribution in this study \\
\midrule
FOEN (Switzerland) & Daily water temperature & Public water-observation service; data are published under the opendata.swiss ``Open use. Must provide the source.'' license. Attribution to FOEN/BAFU and a link to the dataset are required. & Processed annual tables only; daily observations are not redistributed. \\
GKD Bayern (Bavaria) & Daily water temperature & Public official downloads; CC BY 4.0. Attribution: ``Datenquelle: Bayerisches Landesamt f\"ur Umwelt''. & Processed annual tables only; daily observations are not redistributed. \\
USGS NWIS (United States) & Daily water temperature & Public Water Data APIs. U.S. Government work, generally publishable after acquisition; cite USGS Water Data. & Processed annual tables only; daily observations are not redistributed. \\
NASA POWER & Daily air temperature & Public point API. U.S. Government work, generally publishable after acquisition and citation. & Processed annual tables only. \\
LUBW (Baden-W\"urttem\-berg) & Daily water temperature & Restricted. Dataset-specific UIS terms are identified, but their scope and vintage remain pending; redistribution is not established. & Not redistributed. Reviewer access remains pending provider agreement and editorial/legal confirmation. \\
\texttt{gapbench} code & Software & MIT license for original code and project-authored documentation. & Public repository; no credentials required for the code. \\
Processed annual tables & Derived annual quantities & No restricted daily observations; derived from provider records under the terms above. & Included in the public repository; archival DOI pending. \\
\bottomrule
\end{tabular}
\par\smallskip\footnotesize Notes: Code availability does not imply redistribution rights for provider data. The restricted daily observations are not part of the public repository or the Docker image. Source: provider terms and the manuscript's software and data-availability statement; no saved numerical file underlies this table.
\end{table}


\section{Score-bootstrap sensitivity for paired RMSE differences}
Table~\ref{tab:S12} adds score-bootstrap sensitivity intervals for the paired RMSE differences of the main text. The procedure linearizes each RMSE difference through its centered influence function and reweights the 17 connected-group contributions with Rademacher or Webb six-point signs; it never resamples clusters (Kline and Santos, 2012). Its median interval width is 0.88 of the group-resampling width. Every direction that excludes zero under group resampling also excludes zero here: the calendar CDF interval includes zero under both procedures, and the calendar $S_{xx}$, all-year and controlled-gap CDF/$S_{xx}$/all-year contrasts exclude zero under both. Two comparisons are method-dependent: the calendar uniform-exact contrast and the calendar D2$-$classic contrast include zero under group resampling but exclude zero under both score weightings, so neither is promoted to a stable ranking. With 17 unequal clusters neither procedure carries a finite-sample coverage guarantee; these intervals describe comparison variability conditional on the frozen per-archive tables and are not new climate samples.

\begin{table}[p]
\centering
\footnotesize
\caption{Score-bootstrap sensitivity intervals for the paired RMSE differences, alongside the saved group-resampling intervals. The score bootstrap linearizes each RMSE difference through its centered influence function and reweights the 17 connected-group contributions with Rademacher (and Webb six-point) signs, 10{,}000 draws, seed 20261102; clusters are never resampled (Kline and Santos, 2012). Positive estimates mean the first rule or estimator has the larger RMSE. The group-resampling columns reproduce the frozen intervals. Units are $^{\circ}$C per decade.}
\label{tab:S12}
\setlength{\tabcolsep}{3pt}
\resizebox{\textwidth}{!}{%
\begin{tabular}{>{\raggedright\arraybackslash}p{0.125\textwidth}>{\raggedright\arraybackslash}p{0.205\textwidth}r@{\hspace{5pt}}r@{\hspace{5pt}}r@{\hspace{5pt}}r@{\hspace{5pt}}r}
\toprule
Family & Contrast & $n$ & Estimate & Group bootstrap [95\%] & Score Rademacher [95\%] & Score Webb [95\%] \\
\midrule
Calendar & Uniform exact $-$ warm-season & 818 & 0.0219 & $-0.0036$, 0.0369 & 0.0038, 0.0402 & 0.0034, 0.0401 \\
\addlinespace
Calendar & CDF $-$ warm-season & 818 & $-0.0152$ & $-0.0451$, 0.0004 & $-0.0346$, 0.0045 & $-0.0351$, 0.0050 \\
\addlinespace
Calendar & $S_{xx}$ $-$ warm-season & 818 & $-0.0277$ & $-0.0497$, $-0.0169$ & $-0.0420$, $-0.0134$ & $-0.0418$, $-0.0133$ \\
\addlinespace
Calendar & All years $-$ warm-season & 818 & $-0.0462$ & $-0.0706$, $-0.0295$ & $-0.0635$, $-0.0289$ & $-0.0637$, $-0.0289$ \\
\addlinespace
Controlled & Uniform exact $-$ warm-season & 805 & 0.0033 & $-0.0051$, 0.0172 & $-0.0063$, 0.0125 & $-0.0062$, 0.0126 \\
\addlinespace
Controlled & CDF $-$ warm-season & 805 & $-0.0175$ & $-0.0274$, $-0.0052$ & $-0.0282$, $-0.0070$ & $-0.0283$, $-0.0069$ \\
\addlinespace
Controlled & $S_{xx}$ $-$ warm-season & 805 & $-0.0436$ & $-0.0649$, $-0.0172$ & $-0.0660$, $-0.0216$ & $-0.0666$, $-0.0209$ \\
\addlinespace
Controlled & All years $-$ warm-season & 805 & $-0.2762$ & $-0.3574$, $-0.1839$ & $-0.3531$, $-0.2001$ & $-0.3559$, $-0.1976$ \\
\addlinespace
Calendar & D2 $-$ classic & 701 & 0.0104 & $-0.0028$, 0.0178 & 0.0017, 0.0192 & 0.0017, 0.0192 \\
\addlinespace
Controlled & D2 $-$ classic & 738 & 0.0172 & 0.0036, 0.0270 & 0.0064, 0.0282 & 0.0065, 0.0282 \\
\bottomrule
\end{tabular}}
\par\smallskip\footnotesize Notes: the score intervals act on the centered influence function of the RMSE difference and need no cluster resampling; their median width is 0.88 of the group-resampling width (all-year contrasts 0.84--0.88, estimator contrasts 0.85--0.93). Two comparisons are method-dependent: the calendar uniform-exact contrast and the calendar D2$-$classic contrast both include zero under group resampling but exclude it under both score weightings, so neither is promoted to a stable ranking. The calendar CDF interval includes zero under both procedures, and every direction that excludes zero under group resampling also excludes zero under the score procedures. A score interval for the paired mean difference of squared errors is also saved. With 17 unequal clusters neither procedure carries a finite-sample coverage guarantee; this table is a sensitivity check, not a certificate. Source: \path{results/ems_revision_20261102/wild_cluster_bootstrap.csv}, \path{results/ems_revision_20261102/wild_cluster_bootstrap.json}; group-resampling columns from \path{results/wrr_final_revision_20261031/retention_paired_intervals.csv} and \path{results/wrr_final_revision_20261031/estimator_paired_comparisons.csv}.
\end{table}


\end{document}
